\section{Expected electron tails and the radius limits}\label{sec:tails}

Fix any admissible sequence from~\eqref{scales:admissible}, with
arbitrary normalized neutral ground states $\Psi_\ell$ at charges
$Z_\ell$. Use the indices $j_\ell$ and retained events $G_\ell$ of
Proposition~\ref{scales:good}.
Throughout this section, $\sigma_\ell$ is the conditional-integral
version from~\eqref{scales:fields} on the full data space, and
expectations are taken under the original enlarged law.
Thus $\Prob(G_\ell)\to1$, and
Corollary~\ref{prof:uniform} gives uniform convergence of
$\sigma_\ell$ to
\[
 \rho_*(x)=kA_*^{3/2}|x|^{-6}
\]
on each fixed compact annulus, uniformly over the data in $G_\ell$.
We now pass from these conditional densities to expected counts of
raw electron positions. Every fixed sphere has zero expected raw
count, since each one-electron marginal is absolutely continuous.

For $0<a<b<\infty$, write
\[
 \mathcal A(a,b)=\{x:a<|x|<b\},\qquad
 N_\ell(a,b)=\#\{i:a s_\ell<|x_i|<b s_\ell\},
\]
and let
\[
 P_\ell(a,b)=\sum_{i=1}^{Z_\ell}
       \int_{s_\ell\mathcal A(a,b)}\mathcal K_{x_i}(y)\dd y
\]
be the mass of the smoothing packets in the same physical annulus.
The packets and their widths are those of~\eqref{inp:packets} for the
$\ell$th system. Conditional integration and a change of variables give
\begin{equation}\label{tail:conditional-mass}
 \int_{\mathcal A(a,b)}\sigma_\ell(x)\dd x
   =s_\ell^3\E\bigl[P_\ell(a,b)\mid\mathcal F_{j_\ell}\bigr].
\end{equation}

\begin{lemma}[Packet supports]\label{tail:packets}
For each fixed $0<a<b<\infty$, all sufficiently large $\ell$ satisfy
\begin{equation}\label{tail:localization}
 P_\ell(a,b)\le N_\ell(a/2,2b)
\end{equation}
for every raw configuration. Moreover, for every fixed
$0<\eta<\min\{a,(b-a)/2\}$, all sufficiently large $\ell$ satisfy
\begin{equation}\label{tail:sandwich}
 P_\ell(a+\eta,b-\eta)
    \le N_\ell(a,b)
    \le P_\ell(a-\eta,b+\eta)
\end{equation}
for every raw configuration.
\end{lemma}

\begin{proof}
Abbreviate $s=s_\ell$ and $r_0=\eps Z_\ell^{-1/3}$, and set
$\delta_s=c_1s^w\to0$. If a packet centered at $z$ has $|z|\ge r_0$,
then its width obeys $t_z\le\delta_s|z|$. Consequently
\begin{equation}\label{tail:relative-support}
 \mathcal K_z(y)\ne0
 \quad\Longrightarrow\quad
 (1-\delta_s)|z|\le |y|\le(1+\delta_s)|z|.
\end{equation}
This estimate holds also for arbitrarily distant centers. If instead
$|z|<r_0$, its packet is contained in the ball of radius
\[
 r_0+c_1r_0^{1+w}=o(s),
\]
because $r_0/s=\eps(Z_\ell s_\ell^3)^{-1/3}\to0$. Such packets
therefore cannot meet any fixed annulus $s\mathcal A(a,b)$ for large
$\ell$. For the remaining packets, \eqref{tail:relative-support}
shows that a packet meeting this annulus has center in
$s\mathcal A(a/2,2b)$ once $\delta_s$ is small. Each packet has mass
one, proving~\eqref{tail:localization}.

The same inequalities show that every packet meeting the contracted
annulus $s\mathcal A(a+\eta,b-\eta)$ has its center in
$s\mathcal A(a,b)$. Conversely, every packet whose center lies in
$s\mathcal A(a,b)$ is entirely contained in the enlarged annulus
$s\mathcal A(a-\eta,b+\eta)$. For example,
$\delta_s<\eta/b$ suffices for these inclusions once the clamped
packets have been excluded. Summing the mass-one packets proves both
inequalities in~\eqref{tail:sandwich}.
\end{proof}

\begin{lemma}[Expected annular mass]\label{tail:annular}
For every fixed $0<a<b<\infty$,
\begin{align}
 \E\int_{\mathcal A(a,b)}\sigma_\ell(x)\dd x
       &\longrightarrow \int_{\mathcal A(a,b)}\rho_*(x)\dd x,
          \label{tail:smoothed-annulus}\\
 s_\ell^3\E N_\ell(a,b)
       &\longrightarrow \int_{\mathcal A(a,b)}\rho_*(x)\dd x.
          \label{tail:raw-annulus}
\end{align}
\end{lemma}

\begin{proof}
Put $Y_\ell=\int_{\mathcal A(a,b)}\sigma_\ell$ and
$X_\ell=N_\ell(a/2,2b)$. Equations~\eqref{tail:conditional-mass}
and~\eqref{tail:localization}, followed by conditional Jensen, give
\[
 0\le Y_\ell\le s_\ell^3\E[X_\ell\mid\mathcal F_{j_\ell}],
 \qquad
 \|Y_\ell\|_2\le s_\ell^3\|X_\ell\|_2.
\]
The annular count bound~\eqref{inp:counts}, with the common offset
$D_0$ and $s_\ell<1$, therefore yields
\begin{equation}\label{tail:uniform-integrability}
 \|Y_\ell\|_2
 \le C\bigl(1+\sqrt{D_0}\,s_\ell^{7/2}\bigr)\le C'.
\end{equation}
In particular, for the actual complement of the retained event,
\[
 \E\bigl[Y_\ell\1_{G_\ell^c}\bigr]
       \le C'\Prob(G_\ell^c)^{1/2}\longrightarrow0.
\]
No independence of $G_\ell$ from the conditional fields is required.
On $G_\ell$, Corollary~\ref{prof:uniform} implies that $Y_\ell$
converges to $\int_{\mathcal A(a,b)}\rho_*$ uniformly over all retained
data. Since $\Prob(G_\ell)\to1$, this proves
\eqref{tail:smoothed-annulus}.

Taking expectations in~\eqref{tail:conditional-mass} gives
$\E Y_\ell=s_\ell^3\E P_\ell(a,b)$. Apply
\eqref{tail:smoothed-annulus} to the contracted and enlarged annuli
in~\eqref{tail:sandwich}. It follows that
\begin{align*}
 \int_{\mathcal A(a+\eta,b-\eta)}\rho_*
  &\le \liminf_{\ell\to\infty}s_\ell^3\E N_\ell(a,b)\\
  &\le \limsup_{\ell\to\infty}s_\ell^3\E N_\ell(a,b)
   \le \int_{\mathcal A(a-\eta,b+\eta)}\rho_*.
\end{align*}
Letting $\eta\downarrow0$ proves~\eqref{tail:raw-annulus}, because
$\rho_*$ is continuous and integrable on a neighborhood of the fixed
closed annulus.
\end{proof}

\begin{lemma}[Exterior tightness]\label{tail:tightness}
There is a constant $C$, independent of the fixed number $M>0$, such
that every admissible sequence satisfies
\begin{equation}\label{tail:tightness-bound}
 \limsup_{\ell\to\infty}s_\ell^3
      \E\#\{i:|x_i|>Ms_\ell\}\le CM^{-3}.
\end{equation}
\end{lemma}

\begin{proof}
For fixed $M$, eventually $Ms_\ell<1$. Cover the region between
$Ms_\ell$ and $1$ by shells $u<|x|<2u$, with
$u=2^pMs_\ell<1$ and $p=0,1,\ldots$. The annular bound
\eqref{inp:counts}, applied with the fixed shape $(1,2)$, gives
\[
 \E\#\{i:u<|x_i|<2u\}
 \le C\bigl(u^{-3}+\sqrt{D_0u}\bigr)\le C_Du^{-3}
 \qquad(0<u<1).
\]
The expected raw mass outside radius $1$ is bounded by the other
estimate in~\eqref{inp:counts}. Summing the geometric series, and
allowing the last shell to overlap this exterior region, gives
\[
 s_\ell^3\E\#\{i:|x_i|>Ms_\ell\}
 \le C_D M^{-3}\sum_{p=0}^{\infty}2^{-3p}+Cs_\ell^3.
\]
The constants depend only on the fixed count estimates and $D_0$;
they do not depend on $M$. Taking the upper limit proves the claim.
\end{proof}

\begin{theorem}[Expected exterior asymptotic]\label{tail:asymptotic}
Along every admissible sequence, with arbitrary normalized neutral
ground states, for each fixed $a>0$,
\begin{equation}\label{tail:asymptotic-equation}
 s_\ell^3\int_{|y|>as_\ell}\rho_{\Psi_\ell}(y)\dd y
       \longrightarrow \frac{b_{\mathrm{TF}}^3}{a^3}.
\end{equation}
\end{theorem}

\begin{proof}
Fix $M>a$. Lemma~\ref{tail:annular} gives convergence on
$\mathcal A(a,M)$, whereas Lemma~\ref{tail:tightness} bounds the
remaining exterior mass after rescaling by $CM^{-3}$. Thus, with
$T_\ell(r)=\int_{|y|>r}\rho_{\Psi_\ell}(y)\dd y$,
\begin{align*}
 \liminf_{\ell\to\infty}s_\ell^3T_\ell(as_\ell)
   &\ge \int_{\mathcal A(a,M)}\rho_*,\\
 \limsup_{\ell\to\infty}s_\ell^3T_\ell(as_\ell)
   &\le \int_{\mathcal A(a,M)}\rho_*+CM^{-3}.
\end{align*}
Let $M\to\infty$. The limiting density is integrable outside every
positive radius, and
\[
 \int_{|x|>a}\rho_*(x)\dd x
 =\frac{4\pi kA_*^{3/2}}{3a^3}
 =\frac{4A_*}{a^3}
 =\frac{b_{\mathrm{TF}}^3}{a^3},
\]
using $4\pi kA_*^{1/2}=12$ and $4A_*=81\pi^2/2$.
This proves~\eqref{tail:asymptotic-equation}.
\end{proof}

\begin{proof}[Proof of Theorem~\ref{intro:main}]
For each sufficiently large integer $m$, put $s=m^{-1/3}$ and choose
a finite integer
\[
 Z_{\min}(m)\ge
 \max\{m+1,\,Z_{\mathrm{bar}}(m^{-1/3}),\,m^{4/3}\}.
\]
The barrier threshold is uniform over normalized ground states.
Consequently every sequence with $m\to\infty$ and
$Z\ge Z_{\min}(m)$ is admissible: in particular,
$Zs^3\ge m^{1/3}\to\infty$.

Fix $0<a_-<b_{\mathrm{TF}}<a_+$. We claim that, for all sufficiently large
integers $m$, all integers $Z\ge Z_{\min}(m)$, and every normalized
neutral ground state $\Psi$ of charge $Z$,
\begin{equation}\label{tail:strict-bracket}
 \int_{|y|>a_-m^{-1/3}}\rho_\Psi(y)\dd y>m,
 \qquad
 \int_{|y|>a_+m^{-1/3}}\rho_\Psi(y)\dd y<m.
\end{equation}
Otherwise choose strictly increasing violating integers $m_\ell$,
charges $Z_\ell\ge Z_{\min}(m_\ell)$, and corresponding ground states
$\Psi_\ell$. With $s_\ell=m_\ell^{-1/3}$ this is an admissible
sequence. Theorem~\ref{tail:asymptotic}, applied separately at $a_-$
and $a_+$, gives normalized tail limits
$b_{\mathrm{TF}}^3/a_-^3>1$ and $b_{\mathrm{TF}}^3/a_+^3<1$. Hence both
inequalities in~\eqref{tail:strict-bracket} eventually hold on this
sequence, a contradiction. This argument includes arbitrary choices
of the ground states at the selected charges.

The tail mass is nonincreasing in its radius and tends to zero, so
the set defining $R_m(\Psi)$ is nonempty. The first inequality in
\eqref{tail:strict-bracket} excludes every radius at most
$a_-m^{-1/3}$ from that set; the second places $a_+m^{-1/3}$ in the
set. Thus
\begin{equation}\label{tail:radius-bracket}
 a_-\le m^{1/3}R_m(\Psi)\le a_+
\end{equation}
for all the states and charges just specified.

Now fix any family $(\Psi_Z)_{Z\ge1}$ of normalized neutral ground
states. For each fixed sufficiently large integer $m$, removing the
finite prefix $Z<Z_{\min}(m)$ changes neither the inner upper limit
nor the inner lower limit. Both
\[
 m^{1/3}\limsup_{\substack{Z\to\infty\\Z>m}}R_m(\Psi_Z)
 \quad\hbox{and}\quad
 m^{1/3}\liminf_{\substack{Z\to\infty\\Z>m}}R_m(\Psi_Z)
\]
therefore lie between $a_-$ and $a_+$. Finally let $m\to\infty$ and
squeeze with arbitrary $a_-<b_{\mathrm{TF}}<a_+$. This proves both asserted
limits. No convergence in $Z$ at fixed $m$ is used.
\end{proof}
